% Generated draft: requires mathematical and source review.
\documentclass[11pt]{article}
\usepackage[margin=25mm]{geometry}
\usepackage{fontspec}
\setmainfont{Latin Modern Roman}
\usepackage{amsmath,amssymb,mathtools,unicode-math}
\setmathfont{Latin Modern Math}
\usepackage{enumitem,needspace,xurl,hyperref}
\setlength{\parindent}{0pt}
\setlength{\parskip}{5pt}
\setlength{\emergencystretch}{3em}
\newcommand{\N}{\mathbb{N}}
\newcommand{\Z}{\mathbb{Z}}
\newcommand{\Q}{\mathbb{Q}}
\newcommand{\R}{\mathbb{R}}
\newcommand{\C}{\mathbb{C}}
\newcommand{\F}{\mathbb{F}}
\newcommand{\A}{\mathbb{A}}
\newcommand{\PP}{\mathbb{P}}
\newcommand{\E}{\mathbb{E}}
\newcommand{\eps}{\varepsilon}
\newcommand{\dd}{\,\mathrm{d}}
\begin{document}
\section*{UnsolvedMath ID: 30004793}
\textbf{Source problem code:} catalogue entry
\subsection{Definitions and mathematical statement}
The supplied record poses four related research questions rather than a single conjecture with specified hypotheses. Let \(\mathcal{Y}\) denote the family of representations obtained by Yu's construction intended in the source discussion. This notation is provisional: the record does not specify the local field, reductive group, coefficient field, admissible construction data, or equivalence relation on representations. Formally, put
\[
\mathcal{Y}_{\mathrm{MG}}
=\{\pi\in\mathcal{Y}:\pi\text{ is a monomial geometric supercuspidal representation in the source's intended sense}\}.
\]
The record does not supply the definition of the predicate appearing in this set, so this expression identifies the requested classification target without providing an independently applicable criterion.

\begin{enumerate}
\item Determine which members of \(\mathcal{Y}\) belong to \(\mathcal{Y}_{\mathrm{MG}}\), and clarify the extent of this subclass. The original wording asks ``how many'' Yu-constructed representations have this property; no counting convention, measure, or notion of prevalence is specified.
\item Determine the extent to which the proof discussed in the source extends beyond monomial representations. The theorem being proved, its other hypotheses, and the intended larger class of representations are not identified in the supplied record.
\item Determine whether the cyclic-base-change assumption used in that proof can be removed, or replaced by another hypothesis while retaining the intended conclusion. The relevant base-change assertion, extensions, groups, and conclusion are not specified here.
\item Determine under what hypotheses Hecke operators at ramified places establish temperedness at those places by the method discussed in the source. The global setting, automorphic representations, local ramification conditions, Hecke operators, and criterion connecting their behavior to temperedness are not specified.
\end{enumerate}

These questions retain the source's unresolved targets. A fully quantified mathematical statement requires the missing definitions and the precise theorem and method referenced by the source; no universal assertion over unspecified fields, groups, or representations is intended.
\subsection{Short English statement}
Classify the Yu-constructed representations that are monomial geometric supercuspidals, determine how far the source's proof extends beyond monomial representations, investigate removing or replacing its cyclic-base-change assumption, and identify when its ramified Hecke-operator method proves temperedness at ramified places. The supplied record omits essential definitions and hypotheses, so these remain explicitly qualified research questions.
\subsection{Sources}
\begin{itemize}
\item Supplied local UnsolvedMath registry record, permanent ID \(30004793\), source problem code \texttt{catalogue entry}, title: \emph{Geometric Supercuspidals, Base Change, and Ramified Temperedness}. Collection: Oberwolfach Reports Open Problems. The record gives proposed year \(2021\) and names no individual proposer.
\item Supplied report citation: \emph{Automorphic Forms, Geometry and Arithmetic} (2022), Oberwolfach Reports, pp.\ 2089--2156, DOI \texttt{10.4171/owr/2021/39}. Original Oberwolfach report: \url{https://doi.org/10.4171/owr/2021/39}. The record describes these reports as documenting mathematical workshops and their open-problem discussions.
\item Supplied literature reference: Lafforgue, \emph{Ramanujan conjecture for automorphic forms}, \url{https://preprints.ihes.fr/storage/OnRamanujan.pdf}.
\item Supplied literature reference, labeled ``geometric types for supercuspidal representations'': \url{https://arxiv.org/abs/1605.08820}.
\item Literature assessment dated exactly \texttt{2026-08-22}: the supplied curation labels the problem \texttt{partially\_solved} and reports proven automorphic applications of geometric and monomial supercuspidals and temperedness results at unramified places, while describing the extent of monomiality, extensions to other representations, and ramified temperedness as unresolved. This is supplied dated curation, not an independently verified literature assessment. The background explicitly requests independent verification of time-sensitive status and cited claims.
\item The registry labels the statement \texttt{corrected\_verified} and states that its standalone formulation and hypotheses were checked against the official Oberwolfach report. This is a supplied verification claim; the sources and that claim have not been independently checked for this draft.
\end{itemize}
\end{document}
